La telemetría incompleta dificulta el pronóstico de recursos de microservicios y puede modificar la precisión que se obtiene al aumentar la complejidad de los modelos. Se examina cómo la descomposición temporal, los descriptores de disponibilidad y el tamaño del modelo afectan al pronóstico de CPU y memoria en cinco fases. El experimento histórico ChronoGraph utiliza 28 servicios, 8005 instantes, división cronológica 70/10/20, historias de doce pasos, tres semillas y veinte máscaras emparejadas. Las variantes E0--E5 y C0--C3/V2 motivan la transferencia calibrada a 32 servicios adicionales. La ampliación compara quince familias con diez semillas neuronales y veinte máscaras. Los experimentos secundarios varían la ausencia, los horizontes directos y las historias. Los modelos congelados se transfieren a un día de Alibaba 2022 tras seis horas de calibración. Históricamente, E5 reduce el MASE emparejado de memoria en 1.280\% frente a DLinear y aumenta el de CPU en 0.984\%. V2 no conserva toda la mejora en memoria. En los servicios adicionales, A3 usa 15 parámetros y mejora el MASE de memoria en 0.02523 frente a E5, de 74 parámetros. Las diez medias por semilla favorecen A3 (prueba exacta de signos con $p=0.0078125$ ajustada por Holm), con mejora en solo 10 de 32 servicios. A3 no mejora de forma general frente a DLinear. PatchTST pequeño y NLinear obtienen las menores medias principales, y persistencia lidera el caso externo. A3 registra medianas CPU/CUDA de 0.050/0.141 ms con entradas residentes para un origen y 32 servicios. Los compromisos observados dependen de las pérdidas sintéticas, los períodos compartidos y la corta evaluación externa.
